\chapter{Économétrie haute fréquence}\label{ch:qm:high-frequency-econometrics}

Mesurée à partir des variations à la seconde du dernier cours échangé, la volatilité annualisée d'une action atteint 60~\%~; mesurée à partir des variations sur cinq minutes, elle vaut 25~\%. La
différence n'est pas celle du marché. Les transactions alternent entre le cours acheteur et le cours vendeur, et les prix sont posés sur une grille au centime~; à une seconde, ces deux
effets pèsent pour l'essentiel dans la variation mesurée, et à cinq minutes presque pas du tout. Le chapitre~18 annonçait que la \omterm{def:qm:volatility-models:rv}{variance réalisée} converge vers la \omterm{def:qm:volatility-models:rv}{variance intégrée} vraie à mesure que
l'échantillonnage devient plus fin~; ce chapitre explique pourquoi, sur des prix réels, elle fait l'inverse, et comment y remédier~: le bruit
et le \omterm{def:qm:high-frequency-econometrics:signature}{signature plot} qui le révèle, l'intervalle d'échantillonnage optimal, des \omterm{def:qm:estimation:estimator}{estimateurs} qui suppriment le bruit au lieu de l'éviter, la corrélation qui
s'annule lorsque deux actifs sont échangés à des moments différents, et les sauts.

Le marché du chapitre est simulé, afin que chaque estimation puisse être confrontée à la vérité~: un log-prix efficient de volatilité annuelle de 25~\%, autour de
36~dollars, avec une transaction par seconde, au cours acheteur ou au cours vendeur d'une grille au centime (le cours acheteur est le prix efficient arrondi au centime inférieur, le cours vendeur
est supérieur d'un centime), sur des journées de 6.5~heures, soit 23\,400 secondes.

\section{Bruit de microstructure et signature plot}

\begin{definition}[Bruit de microstructure, rebond bid--ask]\label{def:qm:high-frequency-econometrics:noise}
Le \emph{bruit de microstructure}\index{bruit de microstructure} est l'écart entre un prix observé et le prix efficient, créé par le processus de négociation~: la
fourchette, la grille des prix, le caractère discret et l'obsolescence des cotations. Le \emph{rebond bid--ask}\index{rebond bid--ask} est sa forme la plus connue~: des transactions successives
au cours acheteur et au cours vendeur font alterner en signe les rendements observés lorsque le prix efficient n'a pas bougé.
\end{definition}

Si le log-prix observé est $Y_i = X_i + u_i$, avec un bruit $u_i$ iid de variance $\omega^2$ indépendant du log-prix efficient $X$, chaque rendement observé porte $u_i - u_{i-1}$, de sorte que
la \omterm{def:qm:volatility-models:rv}{variance réalisée} sur $n$ rendements vérifie
\[
\E[\mathrm{RV}_n] = \mathrm{IV} + 2n\omega^2:
\]
un biais qui croît linéairement avec la fréquence d'échantillonnage et submerge la \omterm{def:qm:volatility-models:rv}{variance intégrée} à haute fréquence. Le bruit rend aussi les rendements consécutifs
corrélés négativement, avec une première autocovariance égale à $-\omega^2$.

\begin{definition}[Estimateur de Roll]\label{def:qm:high-frequency-econometrics:roll}
\emph{estimateur de Roll}\index{estimateur de Roll} (Roll, 1984) de la fourchette bid--ask effective vaut $2\sqrt{-\gamma_1}$, où $\gamma_1$ est la première autocovariance des variations de prix
des transactions~: sous un simple rebond autour d'un prix efficient, les transactions se font au cours médian $\pm c/2$, avec des signes indépendants, et $\gamma_1 = -c^2/4$.
\end{definition}

Sur l'action simulée, dont la fourchette vaut un centime, l'\omterm{def:qm:high-frequency-econometrics:roll}{estimateur de Roll} donne en moyenne 1.14 centimes sur trente jours~: proche, et biaisé à la hausse parce que le prix efficient
déplace aussi le cours acheteur à l'intérieur de la grille.

\begin{definition}[Signature plot]\label{def:qm:high-frequency-econometrics:signature}
Un \emph{signature plot}\index{signature plot} trace la volatilité réalisée moyenne en fonction de l'intervalle d'échantillonnage~: plate
sans bruit, elle monte vers les intervalles les plus courts en présence de bruit (Andersen, Bollerslev, Diebold et Labys, 2000, qui l'ont nommé volatility signature plot).
\end{definition}

La \Cref{fig:qm:high-frequency-econometrics:signature} montre le \omterm{def:qm:high-frequency-econometrics:signature}{signature plot} du chapitre. À partir des transactions, la volatilité réalisée vaut 60.1~\% à une seconde, 35.1~\% à cinq
secondes, 26.1~\% à une minute et 25.4~\% à cinq minutes~; à partir des cotations au point médian, qui ne rebondissent pas mais restent sur la grille, 36.8~\% à une seconde~; à partir du
prix efficient, 25.0~\% à chaque intervalle.

\begin{omfigure}
\begin{tikzpicture}
\begin{semilogxaxis}[width=10cm,height=5cm,xlabel={intervalle d'échantillonnage (secondes)},ylabel={volatilité annualisée (\%)},
  tick label style={font=\small},label style={font=\small},xmin=1,xmax=1800,ymin=20,ymax=65,grid=major,grid style={gray!25},
  xtick={1,10,60,300,1800},xticklabels={1,10,60,300,1800},
  legend style={font=\footnotesize,at={(0.97,0.97)},anchor=north east,fill=white,draw=none},table/col sep=comma]
\addplot[omThm,very thick,mark=*,mark size=1.3pt] table[x=seconds,y=trade]{figdata/methods/21-high-frequency-econometrics/signature.csv};
\addplot[black,thick,mark=square*,mark size=1.3pt] table[x=seconds,y=mid]{figdata/methods/21-high-frequency-econometrics/signature.csv};
\addplot[omDef,very thick,dashed] table[x=seconds,y=efficient]{figdata/methods/21-high-frequency-econometrics/signature.csv};
\legend{prix des transactions,cotations au point médian,prix efficient}
\end{semilogxaxis}
\end{tikzpicture}

\omcaption{\omterm{def:qm:high-frequency-econometrics:signature}{Signature plot} de l'action simulée (volatilité efficiente de 25~\%, grille au centime autour de 36~dollars, une transaction par seconde au cours acheteur ou au cours vendeur), moyennée
sur trente jours~: volatilité réalisée en fonction de l'intervalle d'échantillonnage. Le rebond et la grille ajoutent 35 points à une seconde et presque rien au-delà de deux minutes. Données~:
le tutoriel du chapitre, avec graine fixée.}\label{fig:qm:high-frequency-econometrics:signature}
\end{omfigure}

La variance du bruit peut être estimée à partir de la plus haute fréquence elle-même, $\hat\omega^2 = \mathrm{RV}_n/(2n)$ (qui compte aussi $\mathrm{IV}/2n$, négligeable ici)~: 1.75
points de base d'écart-type, contre 1.60 pour la valeur vraie. Échantillonner moins souvent revient à échanger du biais contre de la variance, et ce compromis admet un optimum.

\begin{proposition}[Échantillonnage optimal sous bruit iid]\label{prop:qm:high-frequency-econometrics:optimal}
Avec un bruit iid de variance $\omega^2$ et une volatilité constante, l'erreur quadratique moyenne de $\mathrm{RV}_n$ vaut approximativement $2\,\mathrm{IQ}/n + (2n\omega^2)^2$, où $\mathrm{IQ} =
\int\sigma^4$ est la quarticité intégrée~; elle est minimale en $n^* = (\mathrm{IQ}/(4\omega^4))^{1/3}$ (Bandi et Russell, 2008).
\end{proposition}

\begin{proof}
L'erreur de discrétisation de la \omterm{def:qm:volatility-models:rv}{variance réalisée} a pour variance $2\,\mathrm{IQ}/n$ (celle du chapitre~18, $\sqrt{2/m}$ avec volatilité constante), et le bruit ajoute le biais $2n\omega^2$
(à quoi s'ajoutent des termes de variance d'ordre inférieur pour un $n$ modéré). En annulant la dérivée $-2\,\mathrm{IQ}/n^2 + 8n\omega^4$, on obtient $n^{*3} = \mathrm{IQ}/(4\omega^4)$.
\end{proof}

Avec le bruit estimé, le rapport bruit sur signal $\omega^2/\mathrm{IV}$ vaut $1.24 \times 10^{-4}$ et $n^* = 254$ rendements par jour, soit un toutes les 92 secondes (82 avec le bruit
vrai). La simulation confirme~: la racine de l'erreur quadratique moyenne relative de la \omterm{def:qm:volatility-models:rv}{variance réalisée} quotidienne vaut 4.80 à une seconde, 0.18 à 30 secondes, 0.11 à 60, 0.12 à 120
et 0.14 à cinq minutes (voir la \cref{fig:qm:high-frequency-econometrics:mse}). Un échantillonnage à la seconde surestime la variance d'un facteur 5.8~; l'optimum évite ce biais au
prix d'utiliser seulement une fraction des données.

\begin{omfigure}
\begin{tikzpicture}
\begin{loglogaxis}[width=10cm,height=5cm,xlabel={intervalle d'échantillonnage (secondes)},ylabel={RMSE de RV / IV},
  tick label style={font=\small},label style={font=\small},xmin=1,xmax=1800,ymin=0.08,ymax=6,grid=major,grid style={gray!25},
  xtick={1,10,60,300,1800},xticklabels={1,10,60,300,1800},ytick={0.1,0.2,0.5,1,2,5},yticklabels={0.1,0.2,0.5,1,2,5},table/col sep=comma]
\addplot[omDef,very thick,mark=*,mark size=1.5pt] table[x=seconds,y=rmse]{figdata/methods/21-high-frequency-econometrics/mse.csv};
\addplot[omThm,thick,dashed,sharp plot] coordinates {(92,0.08) (92,6)};
\end{loglogaxis}
\end{tikzpicture}

\omcaption{Racine de l'erreur quadratique moyenne relative de la \omterm{def:qm:volatility-models:rv}{variance réalisée} quotidienne à partir des prix des transactions, en fonction de l'intervalle d'échantillonnage, sur soixante jours simulés. La ligne pointillée
représente l'optimum de Bandi--Russell, de 92 secondes, calculé à partir de la variance du bruit estimée. Données~: le tutoriel du chapitre, avec graine fixée.}\label{fig:qm:high-frequency-econometrics:mse}
\end{omfigure}

\section{Estimateurs robustes au bruit}

L'échantillonnage épars jette la plus grande part des données. Trois familles en exploitent l'intégralité et suppriment le bruit.

\begin{definition}[Variance réalisée à deux échelles]\label{def:qm:high-frequency-econometrics:tsrv}
La \emph{variance réalisée à deux échelles}\index{variance réalisée à deux échelles} (Zhang, Mykland et Aït-Sahalia, 2005) moyenne les \omterm{def:qm:volatility-models:rv}{variances réalisées} des $K$ grilles éparses décalées
et retranche le biais dû au bruit, estimé sur la grille complète~: $\mathrm{TSRV} = \overline{\mathrm{RV}}^{(K)} - (\bar n/n)\mathrm{RV}^{(\mathrm{all})}$, $\bar n = (n - K + 1)/K$.
\end{definition}

La moyenne sur les $K$ grilles réduit l'erreur de discrétisation de l'\omterm{def:qm:estimation:estimator}{estimateur} épars~; la soustraction élimine son biais résiduel $2\bar n\omega^2$. Il est convergent
sous un bruit iid, au taux lent $n^{-1/6}$.

\begin{definition}[Noyau réalisé, estimateur à pré-moyennage]\label{def:qm:high-frequency-econometrics:kernel}
Un \emph{noyau réalisé}\index{noyau réalisé} (Barndorff-Nielsen, Hansen, Lunde et Shephard, 2008) ajoute des autocovariances pondérées des rendements à haute fréquence à leur
variance, $\gamma_0 + \sum_{h=1}^Hk\bigl(\frac{h-1}H\bigr)(\gamma_h + \gamma_{-h})$, avec une fonction de pondération lisse comme celle de Parzen, pour annuler l'autocovariance négative créée par le bruit.
L'\emph{estimateur à pré-moyennage}\index{estimateur à pré-moyennage} (Jacod, Li, Mykland, Podolskij et Vetter, 2009) moyenne les rendements sur de courtes fenêtres se chevauchant, avec une
fonction de pondération, ce qui réduit le bruit dans chaque fenêtre, puis élève les moyennes au carré, les remet à l'échelle et retranche le biais résiduel.
\end{definition}

Le noyau est la variance de long terme du chapitre~11 appliquée à l'intérieur d'une journée~: le bruit rend les rendements à haute fréquence MA(1), et une somme de type HAC retrouve la variance de la
partie efficiente. Sur trente journées simulées, la volatilité annualisée moyenne issue des transactions vaut 25.4~\% avec la RV à cinq minutes, 25.2~\% avec la RV à deux échelles ($K = 300$), 25.8~\% avec la RV à
l'intervalle optimal, 25.1~\% avec le pré-moyennage (fenêtres de 150 secondes) et 25.1~\% avec le \omterm{def:qm:high-frequency-econometrics:kernel}{noyau réalisé} ($H = 60$). Les écarts-types relatifs d'un jour à l'autre les séparent
(\cref{fig:qm:high-frequency-econometrics:estimators})~: 14.3~\% pour la RV à cinq minutes, 11.9~\% pour la RV à deux échelles, 7.9~\% à l'intervalle optimal, 7.2~\% pour le pré-moyennage et 4.8~\% pour le
noyau.

\begin{omfigure}
\begin{tikzpicture}
\begin{axis}[width=9cm,height=4.8cm,xbar,bar width=9pt,
  xlabel={écart-type relatif entre les jours (\%)},
  ytick={0,1,2,3,4},yticklabels={RV 5 minutes,RV à deux échelles,RV à 92 s,pré-moyennage,noyau réalisé},
  tick label style={font=\small},label style={font=\small},
  xmin=0,xmax=16,ymin=-0.6,ymax=4.6,grid=major,grid style={gray!25},table/col sep=comma]
\addplot[fill=omDef!60,draw=omDef,nodes near coords,nodes near coords style={font=\footnotesize,anchor=west,xshift=1pt}]
  table[x=sd,y=k]{figdata/methods/21-high-frequency-econometrics/estimators.csv};
\end{axis}
\end{tikzpicture}

\omcaption{Précision de cinq \omterm{def:qm:estimation:estimator}{estimateurs} journaliers de la volatilité sur les prix des transactions de l'action simulée, trente jours~: l'écart-type de l'estimation journalière, relatif à la
\omterm{def:qm:volatility-models:rv}{variance intégrée} vraie. Tous sont presque sans biais (volatilités moyennes comprises entre 25.1~\% et 25.8~\% pour une vraie valeur de 25~\%). Données~: le tutoriel du chapitre,
avec graine fixée.}\label{fig:qm:high-frequency-econometrics:estimators}
\end{omfigure}

\section{Observation asynchrone et effet Epps}

\begin{definition}[Transactions non synchrones, effet Epps]\label{def:qm:high-frequency-econometrics:epps}
Les \emph{transactions non synchrones}\index{transactions non synchrones} désignent l'observation de différents actifs à des instants différents et aléatoires. L'\emph{effet
Epps}\index{effet Epps} (Epps, 1979) est la chute qui en résulte de la corrélation réalisée vers zéro lorsque l'intervalle d'échantillonnage se raccourcit~: avec un échantillonnage au dernier tick sur une grille fine,
la plupart des intervalles contiennent une variation de prix d'un actif et aucune de l'autre.
\end{definition}

Deux actions simulées, de corrélation vraie égale à 0.6, échangées à des instants de Poisson espacés en moyenne de 5 et de 15 secondes, le montrent clairement
(\cref{fig:qm:high-frequency-econometrics:epps})~: la corrélation réalisée à partir des prix au dernier tick vaut 0.03 à une seconde, 0.21 à dix secondes, 0.49 à une minute, 0.59 à cinq
minutes et 0.63 à trente minutes (où l'erreur d'échantillonnage due à peu de rendements prend le relais).

\begin{definition}[Estimateur de Hayashi--Yoshida, échantillonnage en temps de rafraîchissement]\label{def:qm:high-frequency-econometrics:hy}
L'\emph{estimateur de Hayashi--Yoshida}\index{estimateur de Hayashi--Yoshida} (Hayashi et Yoshida, 2005) de la covariation somme les produits de toutes les paires de rendements, un par
actif, dont les intervalles d'observation se chevauchent, sans interpolation. L'\emph{échantillonnage en temps de rafraîchissement}\index{echantillonnage@échantillonnage en temps de rafraîchissement} échantillonne les deux actifs à la première date à laquelle chacun a échangé
depuis le précédent instant d'échantillonnage.
\end{definition}

L'\omterm{def:qm:high-frequency-econometrics:hy}{estimateur de Hayashi--Yoshida} est sans biais pour la covariation des prix efficients sans bruit, quels que soient les instants d'observation, car chaque morceau de la covariance
est compté une seule fois. Sur les deux actions, il donne une corrélation moyenne de 0.594 (vraie valeur 0.6), avec un écart-type d'un jour à l'autre de 0.030. L'\omterm{def:qm:high-frequency-econometrics:hy}{échantillonnage en temps de rafraîchissement} seul donne 0.48~:
il synchronise l'échantillonnage, mais les derniers prix qu'il apparie ont encore été observés à des moments différents.

\begin{omfigure}
\begin{tikzpicture}
\begin{semilogxaxis}[width=10cm,height=5cm,xlabel={intervalle d'échantillonnage (secondes)},ylabel={corrélation réalisée},
  tick label style={font=\small},label style={font=\small},xmin=1,xmax=1800,ymin=0,ymax=0.75,grid=major,grid style={gray!25},
  xtick={1,10,60,300,1800},xticklabels={1,10,60,300,1800},
  legend columns=3,legend style={font=\footnotesize,at={(0.5,-0.3)},anchor=north,draw=none,fill=none,
  /tikz/every even column/.append style={column sep=6pt}},table/col sep=comma]
\addplot[omDef,very thick,mark=*,mark size=1.5pt] table[x=seconds,y=corr]{figdata/methods/21-high-frequency-econometrics/epps.csv};
\addplot[omThm,very thick,dashed] table[x=seconds,y=hy]{figdata/methods/21-high-frequency-econometrics/epps.csv};
\addplot[black,thin] table[x=seconds,y=truth]{figdata/methods/21-high-frequency-econometrics/epps.csv};
\legend{dernier tick,Hayashi--Yoshida,corrélation vraie}
\end{semilogxaxis}
\end{tikzpicture}

\omcaption{\omterm{def:qm:high-frequency-econometrics:epps}{Effet Epps}~: corrélation réalisée de deux actions simulées (corrélation vraie 0.6) qui échangent à des instants de Poisson indépendants, en moyenne toutes les 5 et 15 secondes,
en fonction de l'intervalle d'échantillonnage, moyennée sur quarante jours~; et estimation de Hayashi--Yoshida à partir de toutes les transactions (0.594). Données~: le tutoriel du chapitre,
avec graine fixée.}\label{fig:qm:high-frequency-econometrics:epps}
\end{omfigure}

\section{Sauts à haute fréquence}

\begin{definition}[Variation bipuissance]\label{def:qm:high-frequency-econometrics:bipower}
La \emph{variation bipuissance}\index{variation bipuissance} (Barndorff-Nielsen et Shephard, 2004) vaut $\frac\pi2\sum_i|r_i||r_{i-1}|$~: produits de rendements absolus adjacents, pondérés par
$1/\E|Z|^2 = \pi/2$.
\end{definition}

La \omterm{def:qm:volatility-models:rv}{variance réalisée} converge vers la \omterm{def:qm:brownian-motion:qv}{variation quadratique} tout entière, soit la \omterm{def:qm:volatility-models:rv}{variance intégrée} plus la somme des carrés des sauts~; la \omterm{def:qm:high-frequency-econometrics:bipower}{variation bipuissance} converge vers la seule \omterm{def:qm:volatility-models:rv}{variance intégrée},
car un saut n'intervient que par des produits avec un rendement voisin, de nature diffusive, qui s'annulent quand l'échantillonnage devient plus fin. La différence estime la partie due aux sauts. À
cinq minutes, sur quarante journées simulées avec un saut de 1~\% à midi (29~\% de la \omterm{def:qm:brownian-motion:qv}{variation quadratique} de ce jour-là), la \omterm{def:qm:volatility-models:rv}{variance réalisée} vaut en moyenne 1.43 fois la \omterm{def:qm:volatility-models:rv}{variance intégrée}
et la \omterm{def:qm:high-frequency-econometrics:bipower}{variation bipuissance} 1.16~; sans saut, 1.00 et 1.00. À cinq minutes, le saut continue de contaminer la \omterm{def:qm:high-frequency-econometrics:bipower}{variation bipuissance} par ses deux produits voisins~; le
remède est un échantillonnage plus fin, qui refait apparaître le bruit, de sorte que la détection des sauts et l'estimation robuste au bruit doivent être combinées en pratique.

\section{Tutoriel~: la volatilité qui dépendait de l'horloge}\label{tut:qm:high-frequency-econometrics:clock}

\begin{tutorial}
\textbf{Objectif.} Mesurer la volatilité d'une action simulée à chaque fréquence d'échantillonnage, trouver l'intervalle optimal et supprimer le bruit au lieu de l'éviter.
\textbf{État final~:} les \cref{fig:qm:high-frequency-econometrics:signature,fig:qm:high-frequency-econometrics:mse,fig:qm:high-frequency-econometrics:epps}, l'intervalle optimal de 92
secondes, et l'erreur journalière de 4.8~\% du noyau.
\begin{tutsteps}
\item \textbf{\omterm{def:qm:high-frequency-econometrics:tsrv}{Variance réalisée à deux échelles} et \omterm{def:qm:high-frequency-econometrics:kernel}{noyau réalisé}.}
\omcode{code/firm/hfvol/firm_hfvol.py}{44}{64}{Variance réalisée à deux échelles et noyau réalisé de Parzen.}\label{lst:qm:high-frequency-econometrics:tsrv}
\item \textbf{La covariance de Hayashi--Yoshida} sans interpolation.
\omcode{code/firm/hfvol/firm_hfvol.py}{100}{110}{L'estimateur de Hayashi--Yoshida.}\label{lst:qm:high-frequency-econometrics:hy}
\item \textbf{Exécuter} \texttt{signature()}, \texttt{optimal\_interval()}, \texttt{estimators()}, \texttt{epps()} et \texttt{jumps()} dans \texttt{qm\_hf.py}, puis
\texttt{fig\_hf.py}.
\end{tutsteps}
\textbf{Pour aller plus loin.} Fixer le prix de l'action à 360~dollars (la grille devient dix fois plus fine par rapport au prix) et retracer le \omterm{def:qm:high-frequency-econometrics:signature}{signature plot}~; laisser les transactions arriver en
grappes de Hawkes (chapitre~7) et voir si la largeur de bande du noyau doit changer.
\end{tutorial}

\section{Construction~: les estimateurs à haute fréquence}\label{bld:qm:high-frequency-econometrics:hfvol}

\begin{build}
\bfield{Rôle} Toutes les volatilités et corrélations intrajournalières qu'utilise la petite entreprise pour la cotation et le risque sont calculées ici à partir des ticks, avec le bruit traité explicitement.
\bfield{Interface} \texttt{rv(logp, step)}, \texttt{rv\_subsampled}, \texttt{noise\_variance}, \texttt{tsrv(logp, k)}, \texttt{realised\_kernel(logp, H)},
\texttt{preaveraged(logp, kn)}, \texttt{bipower}, \texttt{optimal\_n(iq, noise\_var)}, \texttt{roll\_spread(price)}, \texttt{hayashi\_yoshida(t1, x1, t2, x2)},
\texttt{refresh\_times(*times)}, \texttt{previous\_tick(t, x, grid)}.
\bfield{Règles} Aucune \omterm{def:qm:volatility-models:rv}{variance réalisée} n'est publiée sans son intervalle d'échantillonnage~; \omterm{def:qm:estimation:estimator}{estimateurs} robustes au bruit par défaut pour tout échantillonnage plus fin qu'une minute~; covariances d'actifs
asynchrones par Hayashi--Yoshida ou par un \omterm{def:qm:estimation:estimator}{estimateur} synchronisé robuste au bruit, jamais par des rendements au dernier tick sur une grille fine.
\bfield{Tests d'acceptation} \texttt{code/firm/hfvol/tests/}~: RV sans bruit égale VI~; le biais du bruit $2n\omega^2$~; deux échelles, noyau et pré-moyennage sans biais sous un bruit iid~;
la \omterm{def:qm:high-frequency-econometrics:bipower}{variation bipuissance} ignore un saut~; Hayashi--Yoshida sans biais sous des transactions de Poisson asynchrones et exact sur un exemple traité à la main~; temps de rafraîchissement et échantillonnage au dernier tick~;
l'\omterm{def:qm:high-frequency-econometrics:roll}{estimateur de Roll} sur un pur rebond~; le $n$ optimal minimise l'erreur approchée.
\bfield{Pour aller plus loin} Le \omterm{def:qm:high-frequency-econometrics:kernel}{noyau réalisé} multivarié~; des tests de saut fondés sur le rapport entre \omterm{def:qm:high-frequency-econometrics:bipower}{variation bipuissance} et \omterm{def:qm:volatility-models:rv}{variance réalisée}~; un bruit qui dépend du signe de la transaction et de la fourchette.
\end{build}

\begin{omsources}
\item T.~G.~Andersen, T.~Bollerslev, F.~X.~Diebold et P.~Labys, «~Great realizations~», \emph{Risk}, mars 2000, p.~105--108.
\item R.~Roll, «~A simple implicit measure of the effective bid-ask spread in an efficient market~», \emph{Journal of Finance} 39, 1984.
\item T.~W.~Epps, «~Comovements in stock prices in the very short run~», \emph{Journal of the American Statistical Association} 74, 1979.
\item L.~Zhang, P.~A.~Mykland et Y.~Aït-Sahalia, «~A tale of two time scales~», \emph{Journal of the American Statistical Association} 100, 2005.
\item O.~E.~Barndorff-Nielsen, P.~R.~Hansen, A.~Lunde et N.~Shephard, «~Designing realized kernels to measure the ex post variation of equity prices in the presence of
noise~», \emph{Econometrica} 76, 2008.
\item J.~Jacod, Y.~Li, P.~A.~Mykland, M.~Podolskij et M.~Vetter, «~Microstructure noise in the continuous case: the pre-averaging approach~», \emph{Stochastic Processes and
their Applications} 119, 2009.
\item T.~Hayashi et N.~Yoshida, «~On covariance estimation of non-synchronously observed diffusion processes~», \emph{Bernoulli} 11, 2005.
\item F.~M.~Bandi et J.~R.~Russell, «~Microstructure noise, realized variance, and optimal sampling~», \emph{Review of Economic Studies} 75, 2008.
\item O.~E.~Barndorff-Nielsen et N.~Shephard, «~Power and bipower variation with stochastic volatility and jumps~», \emph{Journal of Financial Econometrics} 2, 2004.
\end{omsources}

\section{Exercices}

\begin{exercise}[$\star$]\label{exo:qm:high-frequency-econometrics:1}
Une action a une \omterm{def:qm:volatility-models:rv}{variance intégrée} journalière de $(1\%)^2$ et son bruit iid un écart-type de 2 points de base. Quelle est la \omterm{def:qm:volatility-models:rv}{variance réalisée} espérée calculée sur 23\,400 rendements
à une seconde, et la volatilité journalière implicite~?
\end{exercise}

\begin{exercise}[$\star$]\label{exo:qm:high-frequency-econometrics:2}
Les variations des prix de transaction ont une première autocovariance de $-0.0004$ (dollars au carré). Quelle fourchette l'\omterm{def:qm:high-frequency-econometrics:roll}{estimateur de Roll} implique-t-il~?
\end{exercise}

\begin{exercise}[$\star$]\label{exo:qm:high-frequency-econometrics:3}
Pourquoi les cotations médianes présentent-elles moins de bruit que les prix de transaction dans le \omterm{def:qm:high-frequency-econometrics:signature}{signature plot}, mais encore un peu~?
\end{exercise}

\begin{exercise}[$\star\star$]\label{exo:qm:high-frequency-econometrics:4}
Calculer le nombre optimal de rendements et l'intervalle optimal de Bandi--Russell pour une volatilité journalière de 1~\% (constante) et un bruit de 2 points de base, sur 23\,400 secondes.
\end{exercise}

\begin{exercise}[$\star\star$]\label{exo:qm:high-frequency-econometrics:5}
Montrer qu'avec un bruit iid, les rendements haute fréquence observés ont une première autocorrélation $-\omega^2/(\sigma^2\Delta + 2\omega^2)$, et l'évaluer pour l'action du chapitre à une seconde.
\end{exercise}

\begin{exercise}[$\star\star$]\label{exo:qm:high-frequency-econometrics:6}
Deux actifs s'échangent en moyenne toutes les 5 et 15 secondes. Environ quelle part des intervalles d'une seconde contient une transaction pour les deux~?
\end{exercise}

\begin{exercise}[$\star\star\star$]\label{exo:qm:high-frequency-econometrics:7}
\emph{Programmation.} Tracer le \omterm{def:qm:high-frequency-econometrics:signature}{signature plot} de l'action simulée à partir des transactions, des cotations médianes et des prix efficients, puis calculer la racine de l'erreur quadratique moyenne de la \omterm{def:qm:volatility-models:rv}{variance réalisée} pour
chaque intervalle.
\end{exercise}

\begin{exercise}[$\star\star\star$]\label{exo:qm:high-frequency-econometrics:8}
\emph{Repérer l'erreur.} «~Pour estimer la corrélation de deux actions aussi précisément que possible, nous avons utilisé chaque rendement à une seconde de chacune sur une année~; la corrélation vaut 0.12,
donc la paire est inutile pour la couverture.~»
\end{exercise}

\section{Problème~: la volatilité qui dépendait de l'horloge}

\begin{problem}[{Problème du week-end --- échantillonnage d'un prix bruité}]\label{pb:qm:high-frequency-econometrics:1}
Une action simulée a une volatilité efficiente de 25~\% par an, se négocie une fois par seconde au cours acheteur ou au cours vendeur d'une grille au centime autour de 36 dollars, sur des journées de 23\,400 secondes.

\medskip\noindent\textbf{Partie I --- Le \omterm{def:qm:high-frequency-econometrics:signature}{signature plot}.}
\begin{enumerate}
\item Quelle volatilité donnent les rendements de transactions à une seconde, à cinq secondes, à une minute et à cinq minutes~?
\item Et les cotations médianes à une seconde~? Et le prix efficient~?
\item Quel est l'écart-type du bruit, estimé et vrai~?
\item Que dit l'\omterm{def:qm:high-frequency-econometrics:roll}{estimateur de Roll} de la fourchette~?
\item Pourquoi la variance à une seconde est-elle 5.8 fois trop grande~?
\end{enumerate}

\medskip\noindent\textbf{Partie II --- L'intervalle optimal.}
\begin{enumerate}[resume]
\item Quel rapport bruit sur signal les données impliquent-elles~?
\item Quel est le nombre optimal de rendements et l'intervalle optimal de Bandi--Russell~?
\item Comment la racine de l'erreur quadratique moyenne simulée varie-t-elle avec l'intervalle~?
\item Combien coûte l'optimum en données~?
\item Comment l'optimum changerait-il un jour deux fois plus volatil~?
\end{enumerate}

\medskip\noindent\textbf{Partie III --- Au-delà de l'échantillonnage.}
\begin{enumerate}[resume]
\item Que donnent la RV à deux échelles, le pré-moyennage et le \omterm{def:qm:high-frequency-econometrics:kernel}{noyau réalisé}, et avec quelle précision~?
\item Quel est l'\omterm{def:qm:high-frequency-econometrics:epps}{effet Epps} pour deux actions qui s'échangent toutes les 5 et 15 secondes~?
\item Que donnent Hayashi--Yoshida et l'\omterm{def:qm:high-frequency-econometrics:hy}{échantillonnage en temps de rafraîchissement}~?
\item Que fait un saut de 1~\% à la \omterm{def:qm:volatility-models:rv}{variance réalisée} et à la \omterm{def:qm:high-frequency-econometrics:bipower}{variation bipuissance} à cinq minutes~?
\item Pourquoi ne peut-on pas traiter séparément les sauts et le bruit~?
\end{enumerate}

\medskip\noindent\textbf{Partie IV --- Jugement.}
\begin{enumerate}[resume]
\item Quelle volatilité intrajournalière le moteur de cotation d'un teneur de marché doit-il utiliser~?
\item Quelle corrélation sa couverture doit-elle utiliser~?
\item Que faudrait-il vérifier sur des données réelles que la simulation écarte par hypothèse~?
\item Énoncer le \emph{résultat nommé}~: l'intervalle d'échantillonnage optimal et le biais de la \omterm{def:qm:volatility-models:rv}{variance réalisée} à une seconde qu'il permet d'éviter.
\item En une phrase~: pourquoi davantage de données dégradent-elles la \omterm{def:qm:volatility-models:rv}{variance réalisée}~?
\end{enumerate}
\end{problem}

\section{Questions d'entretien}

\begin{interviewq}[$\star$ \iqroles{chercheur, trader}]\label{iq:qm:high-frequency-econometrics:1}
Pourquoi la volatilité mesurée à partir des données au tick est-elle plus élevée que celle obtenue à partir des données à cinq minutes~?
\end{interviewq}

\begin{interviewq}[$\star\star$ \iqroles{chercheur, développeur}]\label{iq:qm:high-frequency-econometrics:2}
Comment estimeriez-vous la fourchette d'une action à partir des seuls prix de transaction~?
\end{interviewq}

\begin{interviewq}[$\star\star$ \iqroles{chercheur}]\label{iq:qm:high-frequency-econometrics:3}
Quelle est la fréquence d'échantillonnage optimale pour la \omterm{def:qm:volatility-models:rv}{variance réalisée}, et de quoi dépend-elle~?
\end{interviewq}

\begin{interviewq}[$\star\star$ \iqroles{chercheur, trader}]\label{iq:qm:high-frequency-econometrics:4}
Votre corrélation réalisée entre deux actions liquides tombe de 0.6 à 0.1 lorsque vous passez de données à cinq minutes à des données à une seconde. Pourquoi, et comment la corrigeriez-vous~?
\end{interviewq}

\begin{interviewq}[$\star\star$ \iqroles{chercheur}]\label{iq:qm:high-frequency-econometrics:5}
Comment séparez-vous les sauts de la volatilité diffusive dans les données intrajournalières~?
\end{interviewq}

\begin{interviewq}[$\star\star\star$ \iqroles{chercheur}]\label{iq:qm:high-frequency-econometrics:6}
Expliquez pourquoi un \omterm{def:qm:high-frequency-econometrics:kernel}{noyau réalisé} élimine le \omterm{def:qm:high-frequency-econometrics:noise}{bruit de microstructure}, en termes d'autocovariances.
\end{interviewq}
